% General projection/resolvent mathematics

许多有效理论都从同一个线性代数问题开始：完整 Hilbert 空间可分成希望显式保留的目标子空间 $P\mathcal H$ 与被消去的互补子空间 $Q\mathcal H$。定义正交投影
\begin{equation*}
 P^2=P=P^\dagger,
 \qquad
 Q\equiv I-P,
 \qquad
 PQ=QP=0.
\end{equation*}
完整 Hamilton 算符记为 $H$，单位为 J；复能量参数 $z$ 也以 J 为单位，并取在 $QHQ$ 的预解集中，使 $(z-QHQ)^{-1}$ 存在。

把 $z-H$ 按 $P\oplus Q$ 分块，先写成
\begin{equation}
 \boxed{
 z-H=
 \begin{pmatrix}
 z-PHP & -PHQ\\
 -QHP & z-QHQ
 \end{pmatrix}.}
 \label{eq:feshbach-block-matrix-dp68}
\end{equation}
对右下块作 Schur 补，目标子空间中的精确预解式为
\begin{equation}
 \boxed{
 P(z-H)^{-1}P=
 \left[z-PHP-PHQ(z-QHQ)^{-1}QHP\right]^{-1}.}
 \label{eq:feshbach-schur-proof-r10}
\end{equation}
这条式子没有假设 $P$--$Q$ 耦合很弱；它只是分块逆矩阵恒等式。由右端自然定义能量依赖的投影自能和精确有效 Hamilton 算符
\begin{equation}
 \boxed{
 H_{\rm eff}(z)
 \equiv PHP+\Sigma_P(z),
 \qquad
 \Sigma_P(z)
 \equiv PHQ(z-QHQ)^{-1}QHP.}
 \label{eq:feshbach_schur_exact_positive_projection}
\end{equation}
$H_{\rm eff}$ 与 $\Sigma_P$ 的单位均为 J。$\Sigma_P(z)$ 记录从 $P$ 子空间跃迁到 $Q$、在 $Q$ 中传播、再返回 $P$ 的全部线性虚过程；其能量依赖来自被消去子空间的预解式，不能在没有额外尺度条件时删去。

若目标能窗与 $QHQ$ 的谱之间存在有限能隙，能量依赖的精确预解式可在该谱隔离控制下展开为静态分块对角化。把
\begin{equation*}
 H=H_{\rm d}+H_{\rm od},
 \qquad
 H_{\rm d}=PHP+QHQ,
 \qquad
 H_{\rm od}=PHQ+QHP
\end{equation*}
分成分块对角与分块非对角部分，并取反厄米、纯非对角生成元 $S$，使 $\widetilde H=e^SHe^{-S}$。条件 $S^\dagger=-S$ 保证 $e^S$ 幺正，因此这一分块对角化不会改变内积或 Hamilton 算符的厄米性。为消去一阶非对角块，$S$ 必须满足 Sylvester 方程
\begin{equation}
 \boxed{H_{\rm od}+[S,H_{\rm d}]=0.}
 \label{eq:sw-sylvester-r15}
\end{equation}
Sylvester 方程固定 $S$ 的一阶非对角部分；同一个 $S$ 再通过 BCH 的二阶对易子产生目标子空间内的能级位移与诱导耦合。把 Baker--Campbell--Hausdorff 级数保留到 $S^2$，并使用\eqnrefs{eq:sw-sylvester-r15} 消去一阶非对角项后，得到
\begin{equation}
 \boxed{
 \widetilde H
 =H_{\rm d}+\frac12[S,H_{\rm od}]
 +\mathcal O(S^3H).}
 \label{eq:sw-bch-second-r15}
\end{equation}
$1/2$ 来自一阶对易子 $[S,H_{\rm od}]$ 与二阶嵌套对易子 $[S,[S,H_{\rm d}]]/2$ 在使用 Sylvester 方程后的系数组合。

进一步令 $|p\rangle$、$|p'\rangle$ 是 $P$ 子空间中的 $H_{\rm d}$ 本征态，$|q\rangle$ 是 $Q$ 子空间本征态，分别满足能量 $E_p,E_{p'},E_q$。二阶静态 Schrieffer--Wolff 有效 Hamilton 算符的矩阵元为
\begin{equation}
 \boxed{
 \begin{aligned}
 \langle p|H_{\rm eff}^{\rm SW}|p'\rangle
 ={}&E_p\delta_{pp'}\\
 &+\frac12\sum_{q\in Q}H_{pq}H_{qp'}
 \left[
 \frac1{E_p-E_q}
 +\frac1{E_{p'}-E_q}
 \right]
 +\mathcal O(\epsilon_{\rm SW}^3\Delta_Q).
 \end{aligned}}
 \label{eq:sw_arbitrary_target_subspace_matrix_elements}
\end{equation}
其中 $H_{pq}=\langle p|H|q\rangle$，$\Delta_Q$ 是目标能窗到 $QHQ$ 谱的最小能隙，
$\epsilon_{\rm SW}\equiv\|PHQ\|/\Delta_Q$ 是无量纲混合参数。只有 $\epsilon_{\rm SW}\ll1$ 时，上式的二阶截断才受控；否则应回到\eqnrefs{eq:feshbach_schur_exact_positive_projection} 的精确能量依赖形式。
